\section{引言}

部分子分布函数（PDFs）$f(x)$ \cite{Feynman:1973xc}
与光锥 $z^2=0$ 上双定域算符的矩阵元素相联系，
这使得在欧氏空间中表述的
格点规范理论里
无法直接计算这些函数。
通常的出路是计算它们的矩。
然而，最近 X.~Ji \cite{Ji:2013dva} 提出了一种方法，
允许把 PDFs 作为 $x$ 的函数来计算。
为此，他建议采用纯类空的
间隔 $z=(0,0,0,z_3)$。
于是人们处理的是准部分子分布（quasi-PDFs）$Q(y,p_3)$，它描述
强子动量 $p_3$ 分量的分配，
并在 \mbox{$p_3 \to \infty$} 极限下
趋于 PDFs $f(y)$。
同样的方法也可用于分布振幅（DAs）。
quasi-PDFs 的格点计算结果
报道于
文献 \cite{Lin:2014zya,Chen:2016utp,Alexandrou:2015rja}，
介子 quasi-DA 的结果见文献 \cite{Zhang:2017bzy}。

在我们近期的论文 \cite{Radyushkin:2016hsy,Radyushkin:2017gjd} 中，
我们利用虚度分布函数（virtuality distribution functions）
\cite{Radyushkin:2014vla,Radyushkin:2015gpa} 的形式体系，
研究了 quasi-PDFs 与 quasi-DAs 的非微扰 $p_3$ 演化。
我们发现 quasi-PDFs 可以从横动量依赖分布（TMDs）
${\cal F} (x, k_\perp^2)$ 得到。我们
用简单的 TMD 模型构建了
quasi-PDFs 非微扰演化的模型。
我们的结果与格点计算得到的
\mbox{$p_3$ 演化} 图像在定性上一致。

在本文中，我们的第一个目标
是建立这样一幅图像：
quasi-PDFs 是 PDFs 与
强子静止系中部分子原始
动量分布的混合体。
作为中间步骤，我们证明 TMDs 与 quasi-PDFs 之间的联系
\cite{Radyushkin:2016hsy}
仅仅是 Lorentz 不变性的一个直接推论。随后我们证明，
当强子运动时，部分子的 $k_3$ 动量
来自两个源头。强子整体的
运动给出
$xp_3$ 部分，它由 TMD ${\cal F} (x, \kappa^2)$ 对其第一个宗量 $x$
的依赖性支配。
其余部分 $k_3-xp_3$
则由 TMD 对其第二个宗量 $\kappa^2$ 的依赖性支配，
后者描述了静止系中的原始
动量分布。quasi-PDFs 的卷积本质
导致其 \mbox{$p_3$ 演化} 的图像相当复杂，
要足够接近 PDF 极限就需要相当大的动量 $p_3 \sim 3$ GeV。

因此，我们的第二个目标是提出一种从格点提取 PDFs 的替代方案。
为此，我们引入
{\it 伪部分子分布（pseudo-PDFs）} ${\cal P} (x, z_3^2) $，
它把光锥 PDFs $f(x)$ 推广到类空间隔，如
$z=(0,0,0,z_3)$。pseudo-PDFs 是
\mbox{{\it Ioffe 时间} \cite{Ioffe:1969kf}} {\it 分布} \cite{Braun:1994jq}
${\cal M} (\nu, z_3^2)$ 的 Fourier 变换；
后者基本上由形如 $
\langle p | \phi(0) \phi (z)|p \rangle $
的一般矩阵元素给出，
写成 $\nu = p_3 z_3$ 和 $z_3^2$ 的函数。
与 quasi-PDFs 不同，pseudo-PDFs 对所有 $z_3^2$
都具有 ``正则'' 的 $-1 \leq x \leq 1$ 支撑。
当 $z_3\to 0$ 时它们趋于 PDFs，
并且在此极限下表现出通常的微扰演化，
其中 $1/z_3$ 充当演化参数。
最后，我们讨论如何利用 pseudo-PDFs 的这些性质
从格点上提取 PDFs。
